\section{Обзор литературы} 
\subsection{Персонализация обучения}
\label{sec:personalization}

Учащиеся, готовящиеся к ЕГЭ по математике, существенно различаются по исходному уровню знаний, темпу усвоения и характерным ошибкам: одни испытывают трудности с алгеброй и функциями, другие систематически ошибаются в планиметрии, третьи не справляются с задачами на параметры. Универсальный конспект, рассчитанный на «среднего» ученика, не учитывает эту неоднородность. Для сильного ученика такой конспект окажется избыточным: он будет тратить время на повторение уже освоенного материала. Для слабого, напротив, он может быть недостаточным, поскольку пропускает промежуточные шаги рассуждения. Отсюда возникает задача персонализации: учебный материал должен адаптироваться к текущему состоянию знаний конкретного школьника.

Педагогическое обоснование такого подхода представлено в работе Блума~\cite{Bloom1984}, показавшего, что индивидуальное тьюторство повышает результаты обучения на два стандартных отклонения по сравнению с групповым форматом. Интеллектуальные обучающие системы (ИОС) пытаются воспроизвести этот эффект программными средствами. Мета-анализ VanLehn~\cite{VanLehn2011} показал, что ИОС с пошаговой обратной связью приближаются по эффективности к человеческому тьюторству, а Steenbergen-Hu и Cooper~\cite{SteenbergenHu2013} подтвердили положительный эффект ИОС на успеваемость именно по математике в школьном сегменте (K--12).

Вместе с тем эффективность персонализации критически зависит от качества диагностики. Если система неточно оценивает, какие темы ученик уже освоил, адаптация содержания может не только не помочь, но и навредить: например, пропустить критически важное объяснение или, наоборот, перегрузить ученика уже знакомым материалом. Обзор адаптивных платформ~\cite{Tan2025} систематизирует основные барьеры: недостаточную интерпретируемость моделей ученика, сложность интеграции в образовательный процесс и риск подмены обучения выдачей готовых ответов. Последний аспект особенно актуален для систем на основе генеративных моделей.

Формализованную модель компетенций ученика обеспечивает семейство методов отслеживания знаний (Knowledge Tracing, KT). Классический подход --- байесовское отслеживание знаний (Bayesian Knowledge Tracing, BKT)~\cite{CorbettAnderson1994} --- моделирует вероятность освоения навыка как скрытую переменную марковской модели. BKT прост в интерпретации, однако предполагает независимость навыков. Нейросетевые методы, прежде всего глубокое отслеживание знаний (Deep Knowledge Tracing, DKT)~\cite{Piech2015}, позволяют моделировать зависимости между навыками, но ценой снижения интерпретируемости. Обзор~\cite{KTSurvey2022} систематизирует современные подходы к KT, а Lu et al.~\cite{Lu2024AKT} развивают направление, в котором акцент смещается с предсказания ответа на диагностику причин затруднений, что ближе к задаче генерации адресного конспекта.

В данной работе отслеживание знаний выполняет конкретную инженерную функцию: превращает историю взаимодействий пользователя с платформой в вектор компетенций по темам ЕГЭ. Пространство тем основано на кодификаторе ФИПИ~\cite{FIPI2025Demo}, что обеспечивает привязку модели к реальной структуре экзамена. Такая привязка принципиальна: без неё вектор компетенций оставался бы абстрактной конструкцией, не соотнесённой с тем, что именно проверяется на экзамене. Полученный вектор становится управляющим сигналом для генерации конспекта: определяет, какие темы включить, какой уровень детализации выбрать, где добавить примеры и визуальные пояснения. В следующем разделе рассматриваются методы генерации учебного контента, использующие такой вектор в качестве входного параметра.

% --------------------------

\subsection{Генерация учебных материалов}
\label{sec:generation}

Вектор компетенций, описанный в предыдущем разделе, становится полезным только при наличии механизма, который преобразует его в учебный материал. В данном разделе рассматриваются три компонента такого механизма: генерация текста на основе больших языковых моделей (Large Language Model, LLM), обеспечение фактической точности через генерацию с извлечением информации (Retrieval-Augmented Generation, RAG) и автоматическое создание визуальных материалов.

Обзор~\cite{LLM4EduSurvey2024} систематизирует сценарии применения LLM в образовании, выделяя роли модели (помощник ученика, помощник преподавателя, компонент платформы) и фиксируя основные риски: галлюцинации, фактические ошибки, подмену обучения выдачей готового ответа. Для данной работы принципиально, что конспект отличается от диалога с чат-ботом: это учебный артефакт, согласованный по структуре, терминологии и уровню сложности. Gupta et al.~\cite{Gupta2025BeyondFinal} показали, что при математическом тьюторинге LLM дают правильный финальный ответ в 85{,}5\% случаев, но лишь 56{,}6\% диалогов оказываются полностью корректными педагогически. Это расхождение указывает на необходимость внешнего контроля структуры и содержания.

Такой контроль обеспечивают методы управляемой генерации~\cite{Liang2024Controllable}: конспект строится по фиксированной структуре обязательных блоков (определения, примеры, разбор типичных ошибок, визуальные пояснения), а LLM адаптирует содержание каждого блока под профиль ученика. Фиксированная структура выполняет роль внешнего каркаса: она не ограничивает содержание, но гарантирует, что конспект покрывает все необходимые элементы и сохраняет последовательность изложения. Проблема галлюцинаций~\cite{HallucinationSurvey2023} в учебном контексте особенно опасна: некорректное определение или ошибка в формуле может закрепить ложное знание. В отличие от диалогового сценария, где ученик может переспросить и получить уточнение, конспект воспринимается как авторитетный источник, поэтому требования к фактической точности здесь выше. Это обосновывает привязку генерации к внешнему корпусу проверенных материалов.

RAG ~\cite{Lewis2020RAG} решает эту проблему: модель обращается к внешнему корпусу, что повышает точность и позволяет обновлять базу знаний без переобучения. Обзор~\cite{Li2025RAGEdu} фиксирует специфику образовательного RAG: в отличие от общего случая, здесь требуется курируемый корпус (учебники, методические материалы, задания с разборами), а извлечённые фрагменты должны соответствовать уровню ученика. Henkel et al.~\cite{Henkel2024RAGMath} продемонстрировали улучшение качества ответов на математические вопросы при использовании RAG с опорой на учебный источник. В разрабатываемой системе корпусом служат материалы ФИПИ~\cite{FIPI2025Bank}, что обеспечивает соответствие школьной программе и требованиям экзамена.

Третий компонент пайплайна связан с визуализацией. В математическом обучении графики, чертежи и блок-схемы выполняют содержательную функцию: связывают формулу с визуальным представлением, иллюстрируют пространственные отношения. Для ряда тем ЕГЭ визуальное представление является не дополнением к тексту, а основным способом понимания: задачи на графики функций, на взаимное расположение фигур, на анализ производной по графику требуют умения читать и строить диаграммы. Конспект, лишённый визуальной составляющей, для таких тем оказывается неполноценным. Мета-анализ~\cite{Schoenherr2024} подтверждает положительный эффект визуализаций, но указывает, что он зависит от типа визуализации, её связи с текстом и когнитивной нагрузки.

Модели генерации изображений по тексту (text-to-image) плохо подходят для этой задачи: они нарушают точные геометрические отношения и некорректно синтезируют подписи и формулы~\cite{Bosheah2025}. Wei et al.~\cite{Wei2025TextToDiagram} формулируют задачу генерации диаграмм по тексту как отдельную и подчёркивают, что результат должен быть редактируемым и проверяемым. Для геометрии Wang et al.~\cite{Wang2025MagicGeo} показали, что корректные пространственные отношения требуют формальных ограничений, поскольку свободная генерация систематически их нарушает. Эти результаты обосновывают архитектуру, принятую в данной работе: языковая модель формирует JSON-спецификацию визуального элемента, методы ограниченной генерации (constrained generation)~\cite{Willard2023Outlines} гарантируют валидность выходного JSON, а специализированный движок рендеринга строит изображение в формате масштабируемой векторной графики (Scalable Vector Graphics, SVG), обеспечивая геометрическую корректность и возможность адаптации под уровень ученика.

Каждый из трёх компонентов решает свою задачу, но ни один по отдельности не закрывает потребность в персонализированном конспекте с визуальными материалами. Управляемая генерация обеспечивает структуру, но не гарантирует фактическую точность. RAG повышает точность, но не создаёт визуальных материалов. Структурная генерация диаграмм обеспечивает корректные иллюстрации, но без привязки к профилю ученика они остаются универсальными. Их интеграция в единый пайплайн, связанный с моделью компетенций ученика, составляет архитектурную основу разрабатываемой системы.

% --------------------------

\subsection{Существующие решения и методы оценки качества}
\label{sec:analogs}

Сравнительный анализ аналогов решает две задачи: показать, что рынок не пуст и отдельные компоненты предлагаемой системы уже реализованы в тех или иных продуктах, и одновременно продемонстрировать, что ни одно из существующих решений не объединяет персонализацию на основе аналитики ученика, автоматическую генерацию конспектов и программное создание точных математических визуализаций.

В российском сегменте подготовки к ЕГЭ по математике доминируют онлайн-школы с видеоуроками и практическими заданиями: Фоксфорд~\cite{Foxford2025}, Умскул~\cite{Umschool2025}, Школково~\cite{Shkolkovo2025}. Эти платформы предлагают структурированные курсы и обратную связь от преподавателей, однако не реализуют автоматическую генерацию персонализированного учебного контента. Портал РЕШУ~ЕГЭ~\cite{ReshuEGE2025} предоставляет обширный банк заданий с возможностью тематической фильтрации, но без адаптации под профиль ученика. Платформы с элементами адаптивности, такие как Учи.ру~\cite{UchiRu2025} и Яндекс Учебник~\cite{YandexAITutor2025}, используют машинное обучение для подбора заданий и предлагают AI-тьюторов по математике, однако их функциональность сосредоточена на диалоговом формате, а не на генерации конспектов с визуальными материалами.

Среди международных продуктов наиболее близким по идеологии является Khanmigo~\cite{Khanmigo2025} от Khan Academy, который акцентирует педагогическую рамку (не выдавать готовый ответ, направлять к пониманию). Специализированные математические инструменты, такие как ALEKS~\cite{ALEKS2025} и MATHia~\cite{CarnegieMATHia2025}, реализуют сильную адаптивную аналитику на основе теории пространств знаний и когнитивных моделей, но работают в формате пошаговых заданий, а не генерации учебных материалов.

В Таблице~\ref{tab:analogs} представлено сравнение репрезентативных платформ по четырём признакам, непосредственно связанным с архитектурой разрабатываемой системы.

\begin{table}[htbp]
\centering
\caption{Сравнение существующих решений по ключевым функциональным признакам}
\label{tab:analogs}
\small
\begin{tabular}{lcccc}
\hline
\textbf{Платформа} & \textbf{Персонализация} & \textbf{Генерация} & \textbf{Авто-} & \textbf{Аналитика} \\
 & \textbf{контента} & \textbf{конспектов} & \textbf{визуализации} & \textbf{ученика} \\
\hline
Фоксфорд          & --     & --     & --     & базовая \\
Умскул             & --     & --     & --     & базовая \\
РЕШУ ЕГЭ           & --     & --     & --     & базовая \\
Учи.ру             & +      & --     & --     & +       \\
Яндекс Учебник     & +      & частично & --   & +       \\
Khanmigo           & +      & частично & --   & +       \\
ALEKS              & +      & --     & --     & ++      \\
MATHia             & +      & --     & --     & ++      \\
\textbf{Данная работа} & \textbf{+} & \textbf{+} & \textbf{+} & \textbf{+} \\
\hline
\end{tabular}
\end{table}

Как видно из таблицы, платформы с сильной аналитикой (ALEKS, MATHia) не генерируют учебный контент автоматически, а продукты с элементами генерации (Khanmigo, Яндекс Учебник) работают в диалоговом формате и не создают структурированных визуализаций. Ни одно из рассмотренных решений не реализует полный цикл создания персонализированных конспектов, что определяет нишу данной работы.

Это расхождение не случайно: оно отражает разницу в архитектурных приоритетах. Платформы с глубокой аналитикой исторически строились вокруг адаптивного тестирования и подбора заданий, а не генерации объяснительного контента. Продукты с генеративными возможностями, напротив, появились с распространением LLM и пока сосредоточены на диалоговом взаимодействии, где визуальный контент не является ключевым элементом. Предлагаемая система занимает промежуточную позицию: она использует аналитику ученика для управления генерацией, а генерацию выстраивает как создание законченного учебного артефакта с текстом и визуальными материалами.